\section*{\centering Abstract}

\noindent An LLM agent finishes a task in one of two ways: either the model judges its own work complete, which we call advisory feedback, or a checker computes completion and the model's claim carries no authority, which we call binding feedback. The distinction matters because an LLM is a generator rather than a verifier. When the specification is incomplete, it fills the gap from its priors, grades its work against that reconstructed answer key, and can confidently declare wrong work finished. We call this failure the false-DONE: the model does not know when it does not know. Across three preregistered experiments on frozen buggy-Python repair corpora, 2,862 episodes over seven models from five providers at under one dollar of total compute, we isolate this authority placement as the only varying factor. Binding lifts a weak model by 9.2 points while leaving a strong one unchanged. Across a seven-model capability ladder the gain is a window, peaking mid-ladder at +14.9 points, zero at the ceiling, and negative for the weakest model, which fails by thrashing rather than false claiming. Raising task difficulty by composing two and three simultaneous bugs roughly doubles the window's depth, up to +31.7 points, but never moves its ceiling: the strongest model produces zero false claims at every difficulty. The window is therefore a property of the model, not the model--task pair, and the advisory false-DONE rate mediates the entire effect, which makes the intervention predictable: a short probe that counts false claims tells a designer in advance whether completion gating will pay.

